\section{Linker}
\label{sec:toolchain-linker}

The \texttt{viv32-ld} tool combines one or more relocatable VIV object files
(\texttt{.vo}) into a final flat executable image. The linker assigns final object
addresses, resolves symbols, applies relocation records, and writes the
resulting executable byte stream.

Input object files are processed in the order supplied on the command line. This order
determines the layout of the final executable image and therefore forms part of the
program definition.

Typical usage is:

\begin{verbatim}
viv32-ld startup.vo main.vo library.vo
viv32-ld startup.vo main.vo library.vo -o program.bin
\end{verbatim}

If no output file is specified, the linked binary is written to the current working
directory using the default output filename.

The linker supports the following command-line options:

\begin{description}
    \item[\texttt{-o <file>}] Write the linked executable to the specified path.
    \item[\texttt{--map <file>}] Generate a linker map describing final object placement
    and resolved symbol addresses.
    \item[\texttt{-h}, \texttt{--help}] Display command-line usage information.
    \item[\texttt{-V}, \texttt{--version}] Display linker version information.
\end{description}

\subsection{Input order}

Input file order is authoritative. Object files are placed into the final output image
in the order supplied on the command line.

For example:

\begin{verbatim}
viv32-ld vectors.vo crt0.vo main.vo serial.vo -o program.bin
\end{verbatim}

places \texttt{vectors.vo} first, followed by \texttt{crt0.vo},
\texttt{main.vo}, and \texttt{serial.vo}.

This allows startup code, reset vectors, and interrupt vectors to be placed at the
beginning of the program without requiring a special entry-point convention or linker
script in the initial toolchain.

\subsection{Object placement}

Each input object is assigned a final base address. The first object is placed at the
start of the output image. Each following object is placed after the bytes of the
previous object.

All symbols defined in an object are adjusted by that object's final base address:

\begin{verbatim}
final_symbol_address = object_base + symbol_offset
\end{verbatim}

where \texttt{symbol\_offset} is the object-local offset recorded in the symbol table.

Object files are placed at four-byte aligned output addresses. If an object's bytecode
length is not a multiple of four, the linker advances the next object's base address to
the next four-byte boundary. The padding is not a symbol-bearing region and is not
addressable through the source object.

\subsection{Symbol resolution}

The linker builds a final symbol table from the symbols provided by all input objects.

All symbols are imported into a global namespace. Because of name mangling, only
symbols marked with \texttt{.nomangle} retain their original name, so they can be resolved anywhere
in the linking. Local symbols with their names mangled will only be relevant for their
own sections, but exist in the global table for address calculations.

Multiple definitions of the same non-mangled symbol
are an error.

A relocation that refers to an undefined symbol is an error.

\subsection{Relocation application}

After object placement and symbol resolution, the linker applies each relocation record
to the output byte stream.

The linker does not interpret VIV-32 instructions. It follows the relocation metadata
provided by the assembler. Each relocation specifies:

\begin{itemize}
    \item which byte offset to patch,
    \item which symbol to resolve,
    \item which addend to apply,
    \item whether the value is absolute or relative,
    \item whether the encoded field is signed or unsigned,
    \item how far to shift the value right before insertion,
    \item how wide the encoded field is,
    \item and where the encoded field is located within the patched word.
\end{itemize}

If the computed value does not fit in the specified field width, linking fails.

\subsection{Output}

The linker emits a flat binary image. The output contains only executable bytes and
data bytes. Symbol tables and relocation tables are not included in the final binary.

The resulting file may be loaded directly into a VIV-32 emulator, virtual machine, or
future hardware implementation according to the configured memory map.
